\def\tutdate{15.12.2016}
\ifdefined\compileall \else
\ifdefined\compiletype
	\documentclass[handout]{beamer}
\else
	\documentclass{beamer}
	\def\compiletype{livebeamer}
\fi

\usepackage{templates/beamerthemekitwide}

\usepackage[utf8]{inputenc}
\usepackage[T1]{fontenc}
\usepackage[ngerman]{babel}
\usepackage{listings}
\usepackage{hyperref}
\usepackage{graphicx}

\usepackage{amsmath}
\usepackage{amsthm}
\usepackage{amssymb}
\usepackage{polynom}

\usepackage{ifthen}
\usepackage{adjustbox} % for \adjincludegraphics

\newcommand{\markBlue}[1]{\textcolor{kit-blue100}{#1}}
\newcommand{\markGreen}[1]{\textcolor{kit-green100}{#1}}

\newcommand{\pitem}{\pause\item}
\newcommand{\p}{\pause}

% -- MATH MACROS
\newcommand{\thistheoremname}{}
\newcommand{\G}{\mathbb{Z}}
\newcommand{\B}{\mathbb{B}}
\newcommand{\R}{\mathbb{R}}
\newcommand{\N}{\mathbb{N}}
\newcommand{\Q}{\mathbb{Q}}
\newcommand{\C}{\mathbb{C}}
\newcommand{\Z}{\mathbb{Z}}
\newcommand{\F}{\mathbb{F}}
\newcommand{\mi}{\mathrm{i}}
\renewcommand{\epsilon}{\varepsilon}


\newenvironment<>{taskblock}[1]{%
	\setbeamercolor{block title}{fg=kit-orange15,bg=kit-orange100}
	\setbeamercolor{block body}{fg=black,bg=kit-orange30}%
	\begin{block}#2{#1}}{\end{block}}

\setbeamertemplate{enumerate items}[default]

% Aussagenlogik Symbole
\newcommand{\W}{w}
\renewcommand{\F}{f}

% Kodierung
\newcommand{\frepr}{\textbf{repr}}
\newcommand{\fRepr}{\textbf{Repr}}
\newcommand{\fZkpl}{\textbf{Zkpl}}
\newcommand{\fbin}{\textbf{bin}}
\newcommand{\fdiv}{\textbf{ div }}
\newcommand{\fmod}{\textbf{ mod }}

\title[Grundbegriffe der Informatik]{Grundbegriffe der Informatik\\Tutorium 33}
\subtitle{}
\author{Lukas Bach, lukas.bach@student.kit.edu}
\date{\tutdate}

\institute{}

\titlelogo{lukasbach}

\titleimage{bg}
%\titleimage{bg-advent}


\ifthenelse{\equal{\compiletype}{livebeamer}}
	{
		\def\livebeamermode{1}
	}{}

\ifthenelse{\equal{\compiletype}{print}}
	{
		\def\printmode{1}
	}{}

\setbeamercovered{invisible}

%\usepackage[citestyle=authoryear,bibstyle=numeric,hyperref,backend=biber]{biblatex}
%\addbibresource{templates/example.bib}
%\bibhang1em

\begin{document}
	
\selectlanguage{ngerman}


%title page
\begin{frame}
	\titlepage
\end{frame}

%table of contents
\ifdefined\printmode
	\ifdefined\compileall \else
	\begin{frame}{Gliederung}
		\tableofcontents
	\end{frame}
\fi\fi

\fi

% pop quiz

\section{Prädikatenlogik}
\begin{frame}{Grundlagen zu Prädikatenlogik}
	Prädikatenlogik (PL) \ip \markBlue{erweitert} Aussagenlogik durch Ergänzen von ``Prädikaten''\ip , einer Art von Funktionen, die Wahrheitswerte zurückgeben.
	
	\bp
	
	Alphabet der Prädikatenlogik:
	
	\begin{itemize}
		\pitem $\lnot, \land, \lor, \rightarrow, \leftrightarrow, (, )$, also Alphabet der Aussagenlogik.
		\pitem $\forall$ Allquantor \ip ($\forall x$ heißt ``für alle $x$ gilt...)
		\pitem $\exists$ Existenzquantor \ip ($\exists x$ heißt ``es existiert min. ein $x$... für das gilt...)
		\pitem $x,y,z,x_i \in Var_{PL}$ Variablen
		\pitem $c, d, c_i \in Const_{PL}$ Konstanten
		\pitem $f, g, h, f_i \in Fun_{PL}$ Funktionen %TODO Selligkeit ar(f_i) = anzahl parameter
		\pitem $\objequiv$ Objektgleichheit
		\pitem $,$ Komma
	\end{itemize}
\end{frame}

\begin{frame}{Gliederung der Prädikatenlogik}
	\begin{block}{Terme}
		Ein Term ist ein Element aus der Sprache über $A_{Ter} := \{\markBlue{(}, \markBlue{)}, \markBlue{,}\} \cup Var_{PL} \cup Const_{PL} \cup Fun_{PL}$.
	\end{block}

	\bp
	
	\begin{block}{Atomare Formeln}
		Atomare Formeln sind zum Beispiel
		\begin{itemize}
			\pitem Objektgleichheiten $R_1 \objequiv R_2$
			\pitem Prädikat von Termen $p(t_1, t_2, ...)$
		\end{itemize}
	\end{block}

	\begin{block}{Stelligkeit einer Funktion}
		Die Stelligkeit $ar(f) \in \N_+$ einer Funktion gibt die Anzahl der Parameter von $f$ an.
	\end{block}

	\bp

	Was sind die Stelligkeiten folgender Funktionen: $f(a, b, c), g(a), h(a, b)$? \pause $4,1,2$.
	
	\bp
\end{frame}

\ifdefined\compileall
\else


\ifthenelse{\equal{\compiletype}{print}}
{

\begin{frame}{Informationen}
	
	\begin{columns}
		\begin{column}{0.5\textwidth}
			
			\begin{block}{Zum Tutorium}
				\begin{itemize}
					\item Lukas Bach
					\item Tutorienfolien auf: 
					\begin{itemize}
						\item \url{http://gbi.lukasbach.com}
					\end{itemize}
					\item Tutorium findet statt:
					\begin{itemize}
						\item Donnerstags, 14:00 - 15:30
						\item 50.34 Informatikbau, -107
					\end{itemize}
				\end{itemize}
			\end{block}
			
			\begin{block}{Mehr Material}
				\begin{itemize}
					\item Ehemalige GBI Webseite:
					\begin{itemize}
						\item \url{http://gbi.ira.uka.de}
						\item Altklausuren!
					\end{itemize}
				\end{itemize}
			\end{block}
			
		\end{column}
		\begin{column}{0.5\textwidth}
			
			\begin{block}{Zur Veranstaltung}
				\begin{itemize}
					\item Grundbegriffe der Informatik
					\item Klausurtermin:
					\begin{itemize}
						\item 06.03.2017, 11:00
						\item Zwei Stunden Bearbeitungszeit
						\item 6 ECTS für Informatiker und Informationswirte, 4 ECTS für Mathematiker und Physiker
					\end{itemize}
				\end{itemize}
			\end{block}
			
			\begin{block}{Zum Übungsschein}
				\begin{itemize}
					\item Übungsblatt jede Woche
					\item Ab 50\% insgesamt hat man den Übungsschein
					\item Keine Voraussetzung für die Klausur, aber für das Modul
				\end{itemize}
			\end{block}
			
		\end{column}
	\end{columns}
	
\end{frame}

}{}

\ifdefined\livebeamermode
	\begin{frame}
		\includegraphics[width=\linewidth]{images/thatsall.png}
	\end{frame}
\fi

\end{document}

\fi